%% ---- Apêndice A: treino dos recomendadores ------------------------------------
\section{Treino dos recomendadores}
\label{ap:treino}

Os recomendadores explicados são o objeto do estudo, não a contribuição, mas
as explicações só fazem sentido se eles aprenderam padrões sequenciais. Este
apêndice descreve como foram treinados e verificados.

\subsection{Arquitetura e treino comuns}

Nos dois conjuntos usamos o SASRec~\cite{kang2018sasrec} implementado em
PyTorch, treinado com entropia cruzada sobre todos os itens, como propõem
Klenitskiy e Vasilev~\cite{klenitskiy2023turning}, em vez da amostragem de um
negativo do artigo original: Adam com taxa $10^{-3}$, lotes de 128 usuários,
comprimento máximo 128, uma cabeça de atenção, \emph{embeddings} de saída
compartilhados com os de entrada, precisão fp32 e corte de gradiente em norma
1. A divisão é \emph{leave-one-out} (último item para teste, penúltimo para
validação) e a avaliação usa ranking completo sobre o catálogo, excluindo os
itens que o usuário já consumiu~\cite{krichene2020sampled}. O treino para
quando o NDCG@10 de validação não melhora por um número fixo de épocas.

\subsection{MovieLens-20M}

Usamos as interações com filtro 5-core iterativo (138\,493 usuários, 18\,345
filmes, 19{,}98 milhões de interações) e a configuração de Klenitskiy e
Vasilev para esse conjunto: 2 camadas e dimensão 256, com paciência de 10
épocas. O melhor modelo é o da época \val{train.ml.best}
(\val{train.ml.epochs} épocas no total), com HR@10 de \val{train.ml.test.hr} e
NDCG@10 de \val{train.ml.test.ndcg} no teste, próximo do que aqueles autores
reportam ($0{,}298$ e $0{,}183$, com comprimento 200).

\subsection{Zvuk}

\paragraph{Dados.}
O Zvuk~\cite{shevchenko2024zvuk} é um registro de cerca de 245 milhões de
audições de 1{,}5 milhão de faixas, publicado no Kaggle sob licença CC BY 4.0.
Contamos como audição uma reprodução de pelo menos 30 segundos (cerca de 40\%
das reproduções são pulos de poucos segundos) e mantemos só a primeira audição
de cada par (usuário, faixa), porque a tarefa explicada, como no MovieLens,
recomenda itens ainda não consumidos. Mantemos as 50\,000 faixas ouvidas pelo
maior número de usuários, porque o treino com entropia cruzada e cada consulta
do verificador custam proporcionalmente ao catálogo; aplicamos o filtro 5-core
e sorteamos 100\,000 usuários, o que deixa 13{,}2 milhões de interações (cerca
de 132 por usuário).

\paragraph{Busca de hiperparâmetros.}
Não há uma configuração publicada para este protocolo, então fizemos uma busca
em duas rodadas em torno da configuração do MovieLens, variando os três
parâmetros que mais pesam no SASRec com entropia cruzada: dimensão
$\{128,256\}$, camadas $\{2,3\}$ e \emph{dropout} $\{0{,}1;0{,}2;0{,}3\}$, com
taxa de aprendizado, tamanho de lote e comprimento fixos. Na primeira rodada as
12 configurações treinam por 8 épocas; na segunda, as três melhores pelo
NDCG@10 de validação continuam do ponto em que pararam até convergir (no máximo
40 épocas, paciência de 5). Escolhemos o modelo de maior NDCG@10 de validação
(Tabela~\ref{tab:busca-zvuk}). As configurações vão de
\val{train.zv.search.min} a \val{train.zv.search.max} na primeira rodada, e as
três finalistas terminam a \val{train.zv.final.spread} umas das outras: a busca
chega a um platô, e não a um pico frágil. O modelo escolhido (dimensão 256,
3 camadas, \emph{dropout} $0{,}1$) para na época \val{train.zv.best}, com
NDCG@10 de \val{train.zv.test.ndcg} e HR@10 de \val{train.zv.test.hr} no
teste.

\begin{table}[t]
\caption{Busca de hiperparâmetros do SASRec no Zvuk. Rodada 1: NDCG@10 de
validação após 8 épocas. Rodada 2 (as três melhores): épocas treinadas (melhor
época), NDCG@10 de validação e de teste. Em negrito, o modelo escolhido.}
\label{tab:busca-zvuk}
\small
\setlength{\tabcolsep}{4pt}
\begin{tabular}{@{}rrcccccc@{}}
\toprule
 & & & Rodada 1 & \multicolumn{3}{c}{Rodada 2} \\
\cmidrule(lr){4-4}\cmidrule(l){5-7}
Dim. & Cam. & \emph{Dropout} & Val. & Épocas & Val. & Teste \\
\midrule
\tabela{busca-zvuk}
\bottomrule
\end{tabular}
\end{table}

\paragraph{Números publicados.}
Zhelnin et al.~\cite{zhelnin2025largecatalog} reportam NDCG@10 de $0{,}080$
a $0{,}094$ para o SASRec treinado com variantes da entropia cruzada no Zvuk,
com o catálogo inteiro (centenas de milhares de faixas) e divisão temporal
global. Nosso protocolo usa um catálogo de 50\,000 faixas, só a primeira
audição de cada faixa e \emph{leave-one-out}, escolhas que elevam muito o
NDCG@10; os números não são comparáveis, e não conhecemos trabalho que use o
nosso protocolo.

\subsection{Os modelos aprenderam padrões sequenciais?}

Para um estudo de explicações, o recomendador não precisa ser o melhor
publicado, e sim usar a ordem do histórico. Verificamos isso nos dois modelos
com \val{model.users} usuários de teste sorteados
(Tabela~\ref{tab:modelos}), comparando o SASRec com a popularidade, com as
transições a partir do último item (com que frequência cada item seguiu o
último item do usuário nas sequências de treino) e consigo mesmo com os
históricos embaralhados, o teste de Klenitskiy et
al.~\cite{klenitskiy2026patterns}.

\begin{table}[t]
\caption{Verificação dos recomendadores: NDCG@10 / HR@10 no teste
(\val{model.users} usuários sorteados, ranking completo).}
\label{tab:modelos}
\small
\begin{tabular}{@{}lcc@{}}
\toprule
 & MovieLens-20M & Zvuk \\
\midrule
\tabela{modelos}
\bottomrule
\end{tabular}
\end{table}

Os dois modelos superam as duas linhas de base com folga. Embaralhar o
histórico tira \val{model.ml.drop}\% do NDCG@10 no MovieLens e
\val{model.zv.drop}\% no Zvuk: os dois usam a ordem, e o do Zvuk depende dela
ao menos tanto quanto o do MovieLens, cuja receita vem da literatura. No Zvuk,
porém, as transições do último item já são fortes (o SASRec as supera em
\val{model.zv.margin}\%, contra \val{model.ml.margin}\% no MovieLens): o sinal
desse conjunto está concentrado no item mais recente, o que explica
explicações ainda mais concentradas no fim do histórico e um guloso mais fraco
em todos os tamanhos (Seção~\ref{sec:robustez}).
